%! Inference and Experiments — Unit 3, Lesson 3.5 — Guided Notes & Practice (Answer Key)
% MIRROR of ../notes/main.tex: -key in place of -boxes, \termblank -> \termans, \blank -> \ans, and the
% fourth argument of every \pcell filled. Nothing else differs — keep it that way.
% THE in-class document, in the MAIN IDEAS / NOTES shape (see LESSON_SHAPE.md §1), at the
% COMPACT budget: the vocabulary box, then ONE guidednotes table of FOUR rows (three instruction
% rows + Guided Practice), 14 problems.
%   I DO   : the teacher fills the definition lines and works the FIRST problem of each grid
%            (1, 5, 8).
%   WE DO  : the class works the rest; the last row, Guided Practice, is worked together.
%   YOU DO : nothing on this page — ../homework/ IS the individual practice.
% THE TRAP is problem 4 (an observed difference of 1.0 "is still a difference" — but 23 of 100
% reshuffles reach it, so chance explains it).
% THE CRUX is problem 10, in the last instruction row: "the 500 teens were chosen at random, so
% this proves music while studying lowers GPA" — random SELECTION lets us generalize; only
% random ASSIGNMENT lets us conclude cause (EK VAR-3.E.2–E.4).
%
% THE CONTEXTS (instruction, guided practice, and homework each use a different one)
%  Rows 1–3 — studying with music: 24 volunteers randomly assigned to silence or music, silence
%          recalls 2.5 more words on average; a GIVEN dotplot of 100 random reshuffles
%          (2 of 100 at 2.5 or more; 23 of 100 at 1.0 or more). Row 3: a random sample of 500
%          teens, music-studiers have lower GPAs (observational).
%  Row 4 (GP) — a reading app: 40 ninth graders picked at random from the district, randomly
%          assigned app / no app; 1 of 100 reshuffles reaches the 6-point difference.
%
%   Simulated differences (silence − music), 100 reshuffles:
%   -3:1 -2.5:1 -2:3 -1.5:7 -1:11 -0.5:16 0:22 0.5:16 1:11 1.5:7 2:3 2.5:1 3:1
%
% PAGE LOCKSTEP with notes_key/: \termblank -> \termans, \blank -> \ans, 4th arg of \pcell.
% TABLE MECHANICS: inside a cell `\\` ENDS THE ROW — use \par. Long rows are split into
% sub-rows (a bare `\\` then `& ...`) so the table can break between them.
% NO SKETCH QUESTIONS: the simulation dotplot is given and read; students never run it.
\documentclass[12pt]{article}
\usepackage{apstats-article}
\usepackage{apstats-key}   % pulls in -boxes; do NOT also load -boxes

\begin{document}

\pageheader{Unit 3, Lesson 3.5}{Guided Notes \& Practice --- Answer Key}
% NAMESTRIP: no \namedateperiod here — the name row lives on the cover only.

\vspace{2pt}

\boxguard[16]
\begin{vocabbox}
\small
Fill in each term \textbf{as we name it} in the notes below.
\par
\termans{Statistical inference}{drawing conclusions beyond the individuals we measured.}
\par
\termans{Statistically significant}{a result too large to be explained by chance alone.}
\par
\termans{Scope of inference}{what a study lets us conclude: cause and effect, and about whom.}
\par
\termans{Generalize}{apply the results to the population the individuals were picked from.}
\par
\end{vocabbox}

\small
\begin{guidednotes}
% ── ROW 1: could chance alone do this? ────────────────────────────────────────
\mainidea[Could chance do it?]{Significant} &
\textbf{Statistical inference} uses data to draw conclusions beyond the \ans{individuals} we
measured. In an experiment, even if the treatment does nothing, the luck of the random
\ans{assignment} makes the groups differ a little. A difference so large that it would
\ans{rarely} happen by chance alone is \textbf{statistically significant}.
\\
& \begin{minipage}[c]{0.44\linewidth}\raggedright
24 volunteers were randomly assigned to study a word list in silence or with music, 12 each.
The silence group recalled \textbf{2.5 more words} on average.\par\vspace{3pt}
A computer then reshuffled the 24 scores into two new groups of 12, \textbf{100 times}, as if
music made no difference.
\end{minipage}\hfill
\begin{minipage}[c]{0.54\linewidth}\centering
\begin{tikzpicture}[x=0.78cm, y=0.1cm]
  \draw[thick] (-3.4,0) -- (3.4,0);
  \foreach \x in {-3,-2,-1,0,1,2,3} {
    \draw (\x,0) -- (\x,-1.4); \node[below, font=\tiny] at (\x,-1.2) {$\x$};}
  \foreach \x/\n in {-3/1,-2.5/1,-2/3,-1.5/7,-1/11,-0.5/16,0/22,0.5/16,1/11,1.5/7,2/3,2.5/1,3/1}
    \foreach \k in {1,...,\n} \fill[navy] (\x,\k) circle (1.3pt);
  \draw[redacc, dashed, thick] (2.5,2) -- (2.5,14);
  \node[font=\tiny\bfseries, redacc, above] at (2.5,14) {observed 2.5};
  \node[font=\tiny] at (0,-7) {Simulated difference in means, silence $-$ music (words)};
  \node[font=\scriptsize\bfseries] at (0,26) {100 random reshuffles};
\end{tikzpicture}
\end{minipage}
\\
& \notesprompt{Read the simulation. Each dot is one reshuffle.}
\begin{probgrid}{|Y|Y|}
\pcell{1}{What does one dot stand for?}{1.2cm}{One reshuffle's difference in means if music made no difference.} &
\pcell{2}{How many of the 100 reshuffles gave a difference of 2.5 words or more?}{1.2cm}{2 of the 100 (at 2.5 and 3.0).} \\ \hline
\pcell{3}{Is the observed 2.5-word difference statistically significant? Why?}{1.2cm}{Yes: chance alone gave 2.5 or more only 2 times in 100.} &
\pcell{4}{Suppose the silence group had recalled only 1.0 more word. Significant? Why?}{1.2cm}{No: 23 of 100 reached 1.0, so chance could explain it.} \\ \hline
\end{probgrid}
\\ \hline
% ── ROW 2: the scope of inference ─────────────────────────────────────────────
\mainidea[What can we say?]{Scope of inference} &
The \textbf{scope of inference} says what a study lets us conclude. Random \ans{assignment}
of treatments lets us conclude \textbf{cause and effect}. Random \ans{selection} of the
individuals from a population lets us \textbf{generalize} to that \ans{population}.
\par\vspace{4pt}
\begin{minipage}{\linewidth}\centering\footnotesize
\renewcommand{\arraystretch}{1.35}
\begin{tabular}{|>{\centering\arraybackslash}m{2.7cm}|>{\raggedright\arraybackslash}m{3.6cm}|>{\raggedright\arraybackslash}m{3.6cm}|}
\hline
& \multicolumn{1}{c|}{\cellcolor{sky}\textbf{Selected at random}}
  & \multicolumn{1}{c|}{\cellcolor{sky}\textbf{Not selected at random}} \tabularnewline
\hline
\cellcolor{sky}\textbf{Assigned at random} & Cause? \ans{Yes}\newline Generalize? \ans{Yes}
  & Cause? \ans{Yes}\newline Generalize? \ans{No} \tabularnewline
\hline
\cellcolor{sky}\textbf{Not assigned} & Cause? \ans{No}\newline Generalize? \ans{Yes}
  & Cause? \ans{No}\newline Generalize? \ans{No} \tabularnewline
\hline
\end{tabular}
\end{minipage}
\\
& \notesprompt{Cause? Generalize to whom?}
\begin{probgrid}{|Y|Y|Y|}
\pcell{5}{300 students picked at random from the district are randomly assigned to music or silence.}{2.0cm}{Cause: yes. Generalize: all students in the district.} &
\pcell{6}{200 students picked at random from the school are asked if they study with music, and their GPA.}{2.0cm}{Cause: no. Generalize: all students at the school.} &
\pcell{7}{The 24 volunteers above, randomly assigned to music or silence.}{2.0cm}{Cause: yes, but only for volunteers like these 24.} \\ \hline
\end{probgrid}
\\ \hline
% ── ROW 3: cause or only association — THE CRUX ───────────────────────────────
\mainidea[Association or]{Cause?} &
With no random assignment, a \ans{confounding} variable may explain an association, so we can
say only that the variables are \ans{associated}. A statistically significant difference in an
experiment with random assignment is evidence that the treatment \ans{caused} it.
\par\vspace{3pt}
A random sample of \textbf{500 teens} from across the state: teens who study with music have
lower GPAs, on average, than teens who study in silence.
\notesprompt{No one was assigned to anything.}
\begin{probgrid}{|Y|Y|Y|}
\pcell{8}{Can we say the association holds for all teens in the state? Why?}{2.2cm}{Yes: the 500 were picked at random from all teens in the state.} &
\pcell{9}{Name a variable that could explain the lower GPAs instead of music.}{2.2cm}{Phone use or texting while studying; total study time.} &
\pcell{10}{``The 500 were chosen at random, so this proves music lowers GPA.'' Agree? Why?}{2.2cm}{No: no one was assigned music. Phone use could explain it.} \\ \hline
\end{probgrid}
\\
& \fcolorbox{navy}{sky}{\parbox{\dimexpr\linewidth-2\fboxsep-2\fboxrule\relax}{\centering
\textbf{Random selection $\to$ generalize. Only random \ans{assignment} $\to$ cause and effect.}}}
\\ \hline
% ── ROW 4: guided practice — WE DO, together, fresh context ──────────────────
\mainidea[Guided practice]{The Reading App} &
A district picks \textbf{40 ninth graders at random} from all of its ninth graders and randomly
assigns 20 to use a reading app for 10 minutes a day and 20 to no app. After 8 weeks the app
group's mean reading-score gain is \textbf{6 points higher}. In 100 random reshuffles of the 40
gains (as if the app did nothing), only \textbf{1} gave a difference of 6 points or more.
\notesprompt{We work these together.}
\begin{probgrid}{|Y|Y|}
\pcell{11}{Is the 6-point difference statistically significant? Why?}{2.0cm}{Yes: only 1 of 100 reshuffles reached 6 points, so chance alone is unlikely.} &
\pcell{12}{Can we conclude the app caused the higher gains? Why?}{2.0cm}{Yes: the groups were assigned at random and the difference is significant.} \\ \hline
\pcell{13}{To whom can we generalize the result? Why?}{2.0cm}{All ninth graders in the district: the 40 were picked at random from them.} &
\pcell{14}{A second district picks 200 ninth graders at random; those who already use the app score higher. ``Random sample, so the app works.'' Respond.}{2.0cm}{No: random sample, not random assignment. Students who chose the app may already read more.} \\ \hline
\end{probgrid}
\\ \hline
\end{guidednotes}

% ── THE NOTES END AT GUIDED PRACTICE ─────────────────────────────────────────
% There is deliberately no "you do" here: ../homework/ is the individual practice,
% started in class in the last 8 minutes while the teacher circulates.

\end{document}
